\documentclass[11pt]{article}
\usepackage[margin=1in]{geometry}
\usepackage[T1]{fontenc}
\usepackage{graphicx}
\usepackage{hyperref}
\usepackage{enumitem}
\setlength{\parskip}{0.5em}
\setlength{\parindent}{0pt}

\title{large language model agents and tool use}
\author{arXiv GraphRAG --- phase-6 deep research}
\date{\today}

\begin{document}
\maketitle

\begin{abstract}
The review is excellent and covers the core trajectory well, but lacks depth in the emerging intersection of world models and test-time compute scaling.
\end{abstract}

\section{From Text Generation to Autonomous Action: A Review of Large Language Model Agents and Tool Use}

The rapid advancement of Large Language Models (LLMs) has fundamentally altered the landscape of artificial intelligence, yet these models remain inherently constrained by the static nature of their training data. While highly capable of linguistic reasoning, LLMs often struggle with tasks requiring real-time information retrieval, precise mathematical computation, or direct interaction with external software. To overcome these boundaries, the field has pivoted toward the development of LLM agents—systems that do not merely predict the next token but actively utilize external tools to bridge the gap between internal reasoning and external execution. This transition is critical, as it transforms LLMs from passive knowledge repositories into proactive problem-solving entities capable of navigating the complexities of the digital world.

This review provides a comprehensive overview of this evolution, tracing the development of agentic capabilities through three distinct stages. We begin by examining the \textbf{Origins} of the field, focusing on early methodologies that introduced basic tool-calling and API integration. We then analyze the \textbf{Key Advances} that moved the field beyond simple execution toward sophisticated reasoning loops, where models began to combine thought, action, and observation to handle more complex workflows. Finally, we investigate the \textbf{Current Frontier}, where research is pushing toward multi-agent collaboration, long-horizon autonomous planning, and the integration of agents into highly dynamic, real-world environments.

\subsection{Origins}

The evolution of large language model (LLM) agents is rooted in the transition from fundamental attention-based architectures to complex systems capable of external knowledge retrieval and embodied interaction. The foundational shift began with the Transformer architecture (\texttt{arXiv:1706.03762}), which introduced an encoder-decoder configuration relying solely on attention mechanisms, effectively replacing recurrence and convolutions for tasks such as machine translation (achieving 28.4 BLEU on WMT 2014 English-to-German) [\texttt{arXiv:1706.03762}]. This architecture was subsequently adapted into bidirectional representations by BERT (\texttt{arXiv:1810.04805}), which conditioned on both left and right context to improve performance in question answering and language inference, reaching 93.2 Test F1 on SQuAD v1.1 [\texttt{arXiv:1810.04805}].

While these models established linguistic competence, a significant tension emerged regarding how models store and access knowledge. GPT-3 (\texttt{arXiv:2005.14165}) demonstrated that massive autoregressive models with 175 billion parameters could perform few-shot learning via text interaction without gradient updates [\texttt{arXiv:2005.14165}]. However, this parametric approach faces challenges, including performance struggles on specific datasets and the inherent difficulty of scaling neural parameters to cover all world facts [\texttt{arXiv:2005.14165}, \texttt{arXiv:2002.08909}]. To address this, research pivoted toward retrieval-augmented architectures. REALM (\texttt{arXiv:2002.08909}) introduced a latent knowledge retriever to attend over large document corpora like Wikipedia, outperforming previous methods in Open-QA by 4-16\% [\texttt{arXiv:2002.08909}]. Similarly, RAG (\texttt{arXiv:2005.11401}) combined parametric memory with a dense vector index to assist knowledge-intensive tasks [\texttt{arXiv:2005.11401}]. Despite these advancements, RAG-based systems still struggle with providing decision provenance and updating world knowledge [\texttt{arXiv:2005.11401}].

The complexity of agentic reasoning is further highlighted by the need for multi-hop reasoning, as seen in the HotpotQA dataset, which addresses the failure of existing QA datasets to support complex factoid comparison and explanation [\texttt{arXiv:1809.09600}]. Furthermore, the limitations of LLMs are exposed by MMLU (\texttt{arXiv:2009.03300}), which identifies critical weaknesses such as near-random accuracy on socially important subjects and a frequent inability for models to recognize when they are wrong [\texttt{arXiv:2009.03300}]. Finally, the trajectory of these models moves toward embodied AI, where agents must bridge the gap between abstract reasoning and concrete execution. The ALFWorld framework (\texttt{arXiv:2010.03768}) addresses this by utilizing modular designs like BUTLER to align text-based policies with visually grounded actions in interactive environments [\texttt{arXiv:2010.03768}]. This progression from pure attention mechanisms to retrieval-augmented and embodied systems marks the foundational shift from static language processing to active, tool-using agency.

\subsection{Key Advances}

The evolution of LLM agents is characterized by a transition from simple prompting to complex, multi-step reasoning and sophisticated tool integration. Early foundational work established Chain-of-Thought (CoT) prompting (\texttt{arXiv:2201.11903}) to elicit intermediate reasoning steps, which significantly improved performance on mathematical tasks like GSM8K. However, CoT is limited by its reliance on manual exemplar composition and a lack of logical guarantees, often producing "fluent but illogical" reasoning paths (\texttt{arXiv:2201.11903}). To address these limitations, the Tree of Thoughts (ToT) framework (\texttt{arXiv:2305.10601}) introduces deliberate decision-making by exploring multiple reasoning paths with lookahead and backtracking, a significant advancement over the token-level, left-to-right constraints of standard CoT. Further advancing this, Reasoning via Planning (RAP) (\texttt{arXiv:2305.14992}) utilizes Monte Carlo Tree Search (MCTS) to balance exploration and exploitation, effectively treating the LLM as both a reasoning agent and a world model.

Tool use has seen a parallel trajectory, moving from self-supervised learning to massive API integration. Toolformer (\texttt{arXiv:2302.04761}) pioneered self-supervised training to decide when and how to call APIs, such as calculators or search engines. In contrast, the ART framework (\texttt{arXiv:2303.09014}) automates multi-step reasoning by selecting demonstrations from a task library to generate programs. While programmatic approaches like Program-Aided Language models (PAL) (\texttt{arXiv:2211.10435}) offload computation to symbolic interpreters to prevent arithmetic errors, other methods focus on scale. ToolLLM (\texttt{arXiv:2307.16789}) facilitates mastery over over 16,000 real-world RESTful APIs, while Gorilla (\texttt{arXiv:2305.15334}) employs retrieval-augmented generation to adapt to frequently updated API documentation. A distinct architectural shift is seen in ToolkenGPT (\texttt{arXiv:2305.11554}), which represents tools as learned embeddings ("toolkens"), allowing the model to trigger tool calls through its standard token-generation mechanism.

The integration of reasoning and acting is best exemplified by ReAct (\texttt{arXiv:2210.03629}), which interleaves verbal reasoning traces with task-specific actions to create a closed-loop synergy. While ReAct provides significant advantages in interactive environments, it remains susceptible to repetitive reasoning loops and is highly dependent on the quality of external tools (\texttt{arXiv:2210.03629}). This complexity is further compounded in multi-agent systems. While CAMEL (\texttt{arXiv:2303.17760}) utilizes role-playing for communicative exploration, MetaGPT (\texttt{arXiv:2308.00352}) introduces a more structured assembly line paradigm, using Standardized Operating Procedures (SOPs) to mitigate the cascading hallucinations often found when naively chaining LLMs. Despite these advances, a significant performance gap remains; for instance, in the WebArena benchmark (\texttt{arXiv:2307.13854}), GPT-4 achieved only a 14.41\% success rate compared to 78.24\% for humans, highlighting the ongoing challenge of long-horizon, real-world task completion.

\subsection{Current Frontier}

The frontier of LLM agents is characterized by a transition from static architectural taxonomies to dynamic, reasoning-centric paradigms. Recent literature suggests a shift from viewing agents as mere collections of components (memory, tools, perception) toward a focus on the underlying cognitive processes that drive behavior [\texttt{arXiv:2601.12538}]. This evolution is increasingly framed through the lens of "Agentic Reasoning," which distinguishes between foundational single-agent planning and collective multi-agent coordination [\texttt{arXiv:2601.12538}].

A central theme in current research is the methodological progression of tool use. Rather than treating tool integration as a single task, recent frameworks categorize the evolution of tool use into three optimization paradigms: prompting as a plug-and-play mechanism, supervised tool learning, and reward-driven tool policy learning [\texttt{arXiv:2604.00835}]. This mirrors the broader shift in agent development from in-context orchestration toward post-training optimization via reinforcement learning (RL) [\texttt{arXiv:2601.12538}]. For instance, the ARTIST framework [\texttt{arXiv:2505.01441}] and DeepSeek-R1 [\texttt{arXiv:2501.12948}] demonstrate the power of RL in incentivizing emergent reasoning patterns, such as self-reflection and verification, which are essential for complex, multi-turn tool invocation [\texttt{arXiv:2505.01441}].

The complexity of agentic workflows is also driving new architectural abstractions. While earlier surveys focused on single-agent and multi-agent systems (MAS) [\texttt{arXiv:2402.01680}, \texttt{arXiv:2401.03428}], newer perspectives introduce a "skill" layer—reusable procedural artifacts that coordinate tools and memory to prevent the inefficiency of reasoning from scratch for every low-level action [\texttt{arXiv:2605.07358}]. This is complemented by efforts to move from simple single-tool calls to sophisticated multi-tool orchestration over long trajectories [\texttt{arXiv:2603.22862}]. To address the high inference costs associated with repeated LLM calls in these trajectories, frameworks like AutoTool [\texttt{arXiv:2511.14650}] utilize graph-based structures to exploit tool-usage inertia, reducing costs by up to 30\%.

However, significant challenges remain in the transition from controlled environments to "in the wild" deployment. Current agents struggle with the manual parsing of complex tool documentation [\texttt{arXiv:2405.16533}] and the lack of verifiable reasoning in long-horizon interactions [\texttt{arXiv:2508.17281}]. Furthermore, while some frameworks attempt to formalize the "think-act" structure using POMDP frameworks to factorize internal reasoning from external execution [\texttt{arXiv:2601.12538}], the distinction between in-context methods and post-training adaptation remains occasionally blurred in existing taxonomies [\texttt{arXiv:2601.12538}]. Finally, the field faces a crisis of transparency; recent audits reveal that many agent benchmarks fail to disclose critical details such as inference costs, sampling settings, or the specific versions of evaluators used, hindering the reproducibility of agentic progress [\texttt{arXiv:2605.21404}].

\subsection{Tensions and Open Problems}

A primary tension in the literature concerns the architectural locus of reasoning: whether intelligence is best elicited through internal linguistic patterns or through externalized computational symbolic logic. On one hand, early approaches focused on enhancing the model's inherent probabilistic reasoning through prompting strategies, such as the self-consistency methods described in \texttt{arXiv:2203.11171} to improve Chain-of-Thought outputs. Conversely, a growing body of work argues that language models are fundamentally limited by their internal arithmetic and logical consistency, advocating instead for "program-aided" paradigms. This is exemplified by the distinction between the program-of-thought prompting in \texttt{arXiv:2211.12588}, which seeks to disentangle computation from language, and the more structured execution frameworks proposed in \texttt{arXiv:2211.10435}, which treat the model as a high-level controller for formal code execution.

Another significant debate centers on the nature of agentic autonomy and the necessity of closed-loop feedback. Some frameworks emphasize the importance of high-fidelity, static environments for benchmarking, such as the WebShop interaction tasks in \texttt{arXiv:2207.01206}. However, more recent critiques suggest that such environments fail to capture the complexity of real-world agency. For instance, the reflexive learning approach in \texttt{arXiv:2303.11366} argues for agents that can learn from verbal reinforcement and self-correction, a move toward autonomy that contrasts with the more rigid, plan-and-execute structures seen in earlier interactive planning models like those in \texttt{arXiv:2302.01560}. This tension highlights a fundamental question: should an agent be a sophisticated executor of pre-defined plans, or a self-evolving entity capable of iterative environmental reasoning?

Finally, there is a growing disagreement regarding the scalability and methodology of tool integration. While foundational works like \texttt{arXiv:2210.03629} (ReAct) and \texttt{arXiv:2201.11903} established the synergy between reasoning and acting, newer research suggests these methods are insufficient for the complexity of modern software ecosystems. The limitations of specialized, narrow tool-learning benchmarks, such as API-Bank in \texttt{arXiv:2408.0244}, have been highlighted by the push toward massive-scale orchestration frameworks like ToolLLM in \texttt{arXiv:2307.16789}. This debate shifts the focus from whether a model can use a single tool to whether a model can master the orchestration of thousands of diverse APIs, moving from simple tool-calling to complex, multi-step tool-use workflows.

\subsection{Outlook}

The frontier of the field is rapidly shifting from "prompt-based" interaction toward "agentic orchestration," where the focus is no longer on the model's ability to answer questions, but on its ability to manage long-horizon, multi-tool workflows in dynamic environments. As we move toward 2026, the trajectory suggests a convergence of embodied reasoning and autonomous tool mastery, where the primary challenge lies in developing robust, standardized evaluation schemas that can capture the emergent behaviors of agents operating in open-ended, real-world digital and physical spaces.

\section*{References}
\begin{enumerate}\small
  \item Attention Is All You Need (2017). \texttt{arXiv:1706.03762}
  \item Proximal Policy Optimization Algorithms (2017). \texttt{arXiv:1707.06347}
  \item BERT: Pre-training of Deep Bidirectional Transformers for Language
  Understanding (2018). \texttt{arXiv:1810.04805}
  \item HotpotQA: A Dataset for Diverse, Explainable Multi-hop Question
  Answering (2018). \texttt{arXiv:1809.09600}
  \item Language Models are Few-Shot Learners (2020). \texttt{arXiv:2005.14165}
  \item REALM: Retrieval-Augmented Language Model Pre-Training (2020). \texttt{arXiv:2002.08909}
  \item Retrieval-Augmented Generation for Knowledge-Intensive NLP Tasks (2020). \texttt{arXiv:2005.11401}
  \item Measuring Massive Multitask Language Understanding (2020). \texttt{arXiv:2009.03300}
  \item ALFWorld: Aligning Text and Embodied Environments for Interactive
  Learning (2020). \texttt{arXiv:2010.03768}
  \item Evaluating Large Language Models Trained on Code (2021). \texttt{arXiv:2107.03374}
  \item Training Verifiers to Solve Math Word Problems (2021). \texttt{arXiv:2110.14168}
  \item WebGPT: Browser-assisted question-answering with human feedback (2021). \texttt{arXiv:2112.09332}
  \item Measuring Mathematical Problem Solving With the MATH Dataset (2021). \texttt{arXiv:2103.03874}
  \item LoRA: Low-Rank Adaptation of Large Language Models (2021). \texttt{arXiv:2106.09685}
  \item Chain-of-Thought Prompting Elicits Reasoning in Large Language Models (2022). \texttt{arXiv:2201.11903}
  \item Training language models to follow instructions with human feedback (2022). \texttt{arXiv:2203.02155}
  \item PaLM: Scaling Language Modeling with Pathways (2022). \texttt{arXiv:2204.02311}
  \item ReAct: Synergizing Reasoning and Acting in Language Models (2022). \texttt{arXiv:2210.03629}
  \item Self-Consistency Improves Chain of Thought Reasoning in Language Models (2022). \texttt{arXiv:2203.11171}
  \item Do As I Can, Not As I Say: Grounding Language in Robotic Affordances (2022). \texttt{arXiv:2204.01691}
  \item Large Language Models are Zero-Shot Reasoners (2022). \texttt{arXiv:2205.11916}
  \item Language Models as Zero-Shot Planners: Extracting Actionable Knowledge
  for Embodied Agents (2022). \texttt{arXiv:2201.07207}
  \item Emergent Abilities of Large Language Models (2022). \texttt{arXiv:2206.07682}
  \item Inner Monologue: Embodied Reasoning through Planning with Language
  Models (2022). \texttt{arXiv:2207.05608}
  \item WebShop: Towards Scalable Real-World Web Interaction with Grounded
  Language Agents (2022). \texttt{arXiv:2207.01206}
  \item PAL: Program-aided Language Models (2022). \texttt{arXiv:2211.10435}
  \item Least-to-Most Prompting Enables Complex Reasoning in Large Language
  Models (2022). \texttt{arXiv:2205.10625}
  \item Self-Instruct: Aligning Language Models with Self-Generated Instructions (2022). \texttt{arXiv:2212.10560}
  \item A Generalist Agent (2022). \texttt{arXiv:2205.06175}
  \item TALM: Tool Augmented Language Models (2022). \texttt{arXiv:2205.12255}
  \item ProgPrompt: Generating Situated Robot Task Plans using Large Language
  Models (2022). \texttt{arXiv:2209.11302}
  \item Program of Thoughts Prompting: Disentangling Computation from Reasoning
  for Numerical Reasoning Tasks (2022). \texttt{arXiv:2211.12588}
  \item Code as Policies: Language Model Programs for Embodied Control (2022). \texttt{arXiv:2209.07753}
  \item TPTU: Large Language Model-based AI Agents for Task Planning and Tool Usage (2023). \texttt{arXiv:2308.03427}
  \item Large Language Models as Tool Makers (2023). \texttt{arXiv:2305.17126}
  \item Towards Reasoning in Large Language Models via Multi-Agent Peer Review
  Collaboration (2023). \texttt{arXiv:2311.08152}
  \item Controlling Large Language Model-based Agents for Large-Scale
  Decision-Making: An Actor-Critic Approach (2023). \texttt{arXiv:2311.13884}
  \item ART: Automatic multi-step reasoning and tool-use for large language
  models (2023). \texttt{arXiv:2303.09014}
  \item An In-depth Survey of Large Language Model-based Artificial Intelligence
  Agents (2023). \texttt{arXiv:2309.14365}
  \item A Survey on Large Language Model based Autonomous Agents (2023). \texttt{arXiv:2308.11432}
  \item AgentBench: Evaluating LLMs as Agents (2023). \texttt{arXiv:2308.03688}
  \item GPT-4 Technical Report (2023). \texttt{arXiv:2303.08774}
  \item Toolformer: Language Models Can Teach Themselves to Use Tools (2023). \texttt{arXiv:2302.04761}
  \item LLaMA: Open and Efficient Foundation Language Models (2023). \texttt{arXiv:2302.13971}
  \item Llama 2: Open Foundation and Fine-Tuned Chat Models (2023). \texttt{arXiv:2307.09288}
  \item Reflexion: Language Agents with Verbal Reinforcement Learning (2023). \texttt{arXiv:2303.11366}
  \item Generative Agents: Interactive Simulacra of Human Behavior (2023). \texttt{arXiv:2304.03442}
  \item Sparks of Artificial General Intelligence: Early experiments with GPT-4 (2023). \texttt{arXiv:2303.12712}
  \item Retrieval-Augmented Generation for Large Language Models: A Survey (2023). \texttt{arXiv:2312.10997}
  \item Tree of Thoughts: Deliberate Problem Solving with Large Language Models (2023). \texttt{arXiv:2305.10601}
  \item Self-Refine: Iterative Refinement with Self-Feedback (2023). \texttt{arXiv:2303.17651}
  \item ToolLLM: Facilitating Large Language Models to Master 16000+ Real-world
  APIs (2023). \texttt{arXiv:2307.16789}
  \item HuggingGPT: Solving AI Tasks with ChatGPT and its Friends in Hugging
  Face (2023). \texttt{arXiv:2303.17580}
  \item WebArena: A Realistic Web Environment for Building Autonomous Agents (2023). \texttt{arXiv:2307.13854}
  \item Gorilla: Large Language Model Connected with Massive APIs (2023). \texttt{arXiv:2305.15334}
  \item Voyager: An Open-Ended Embodied Agent with Large Language Models (2023). \texttt{arXiv:2305.16291}
  \item CAMEL: Communicative Agents for "Mind" Exploration of Large Language
  Model Society (2023). \texttt{arXiv:2303.17760}
  \item Judging LLM-as-a-Judge with MT-Bench and Chatbot Arena (2023). \texttt{arXiv:2306.05685}
  \item MetaGPT: Meta Programming for A Multi-Agent Collaborative Framework (2023). \texttt{arXiv:2308.00352}
  \item Self-RAG: Learning to Retrieve, Generate, and Critique through
  Self-Reflection (2023). \texttt{arXiv:2310.11511}
  \item Language Models can Solve Computer Tasks (2023). \texttt{arXiv:2303.17491}
  \item Improving Factuality and Reasoning in Language Models through Multiagent
  Debate (2023). \texttt{arXiv:2305.14325}
  \item ChatDev: Communicative Agents for Software Development (2023). \texttt{arXiv:2307.07924}
  \item Describe, Explain, Plan and Select: Interactive Planning with Large
  Language Models Enables Open-World Multi-Task Agents (2023). \texttt{arXiv:2302.01560}
  \item Mind2Web: Towards a Generalist Agent for the Web (2023). \texttt{arXiv:2306.06070}
  \item Tool Learning with Foundation Models (2023). \texttt{arXiv:2304.08354}
  \item The Rise and Potential of Large Language Model Based Agents: A Survey (2023). \texttt{arXiv:2309.07864}
  \item REPLUG: Retrieval-Augmented Black-Box Language Models (2023). \texttt{arXiv:2301.12652}
  \item Augmented Language Models: a Survey (2023). \texttt{arXiv:2302.07842}
  \item SWE-bench: Can Language Models Resolve Real-World GitHub Issues? (2023). \texttt{arXiv:2310.06770}
  \item AutoGen: Enabling Next-Gen LLM Applications via Multi-Agent Conversation (2023). \texttt{arXiv:2308.08155}
  \item API-Bank: A Comprehensive Benchmark for Tool-Augmented LLMs (2023). \texttt{arXiv:2304.08244}
  \item GAIA: a benchmark for General AI Assistants (2023). \texttt{arXiv:2311.12983}
  \item Encouraging Divergent Thinking in Large Language Models through
  Multi-Agent Debate (2023). \texttt{arXiv:2305.19118}
  \item A Survey of Large Language Models (2023). \texttt{arXiv:2303.18223}
  \item Benchmarking Large Language Models in Retrieval-Augmented Generation (2023). \texttt{arXiv:2309.01431}
  \item Ragas: Automated Evaluation of Retrieval Augmented Generation (2023). \texttt{arXiv:2309.15217}
  \item Code Llama: Open Foundation Models for Code (2023). \texttt{arXiv:2308.12950}
  \item Visual ChatGPT: Talking, Drawing and Editing with Visual Foundation
  Models (2023). \texttt{arXiv:2303.04671}
  \item Chameleon: Plug-and-Play Compositional Reasoning with Large Language
  Models (2023). \texttt{arXiv:2304.09842}
  \item Gemini: A Family of Highly Capable Multimodal Models (2023). \texttt{arXiv:2312.11805}
  \item AgentVerse: Facilitating Multi-Agent Collaboration and Exploring
  Emergent Behaviors (2023). \texttt{arXiv:2308.10848}
  \item Reasoning with Language Model is Planning with World Model (2023). \texttt{arXiv:2305.14992}
  \item PaLM-E: An Embodied Multimodal Language Model (2023). \texttt{arXiv:2303.03378}
  \item ChatEval: Towards Better LLM-based Evaluators through Multi-Agent Debate (2023). \texttt{arXiv:2308.07201}
  \item ToolAlpaca: Generalized Tool Learning for Language Models with 3000
  Simulated Cases (2023). \texttt{arXiv:2306.05301}
  \item TaskMatrix.AI: Completing Tasks by Connecting Foundation Models with
  Millions of APIs (2023). \texttt{arXiv:2303.16434}
  \item ToolkenGPT: Augmenting Frozen Language Models with Massive Tools via
  Tool Embeddings (2023). \texttt{arXiv:2305.11554}
  \item Enhancing Retrieval-Augmented Large Language Models with Iterative
  Retrieval-Generation Synergy (2023). \texttt{arXiv:2305.15294}
  \item LLM+P: Empowering Large Language Models with Optimal Planning
  Proficiency (2023). \texttt{arXiv:2304.11477}
  \item A Real-World WebAgent with Planning, Long Context Understanding, and
  Program Synthesis (2023). \texttt{arXiv:2307.12856}
  \item Large Language Model based Multi-Agents: A Survey of Progress and
  Challenges (2024). \texttt{arXiv:2402.01680}
  \item Tool Learning in the Wild: Empowering Language Models as Automatic Tool
  Agents (2024). \texttt{arXiv:2405.16533}
  \item Exploring Large Language Model based Intelligent Agents: Definitions,
  Methods, and Prospects (2024). \texttt{arXiv:2401.03428}
  \item Exploring Autonomous Agents through the Lens of Large Language Models: A
  Review (2024). \texttt{arXiv:2404.04442}
  \item A Survey on Large Language Model-Based Game Agents (2024). \texttt{arXiv:2404.02039}
  \item Large Language Model-Based Agents for Software Engineering: A Survey (2024). \texttt{arXiv:2409.02977}
  \item Learning to Use Tools via Cooperative and Interactive Agents (2024). \texttt{arXiv:2403.03031}
  \item A Review of Prominent Paradigms for LLM-Based Agents: Tool Use
  (Including RAG), Planning, and Feedback Learning (2024). \texttt{arXiv:2406.05804}
  \item The Llama 3 Herd of Models (2024). \texttt{arXiv:2407.21783}
  \item Evaluating Retrieval Quality in Retrieval-Augmented Generation (2024). \texttt{arXiv:2404.13781}
  \item DeepSeekMath: Pushing the Limits of Mathematical Reasoning in Open
  Language Models (2024). \texttt{arXiv:2402.03300}
  \item Large Language Model Agent: A Survey on Methodology, Applications and
  Challenges (2025). \texttt{arXiv:2503.21460}
  \item From Language to Action: A Review of Large Language Models as Autonomous Agents and Tool Users (2025). \texttt{arXiv:2508.17281}
  \item Fundamentals of Building Autonomous LLM Agents (2025). \texttt{arXiv:2510.09244}
  \item AutoTool: Efficient Tool Selection for Large Language Model Agents (2025). \texttt{arXiv:2511.14650}
  \item Large Language Models for Agent-Based Modelling: Current and possible uses across the modelling cycle (2025). \texttt{arXiv:2507.05723}
  \item Generalizability of Large Language Model-Based Agents: A Comprehensive Survey (2025). \texttt{arXiv:2509.16330}
  \item Agentic Reasoning and Tool Integration for LLMs via Reinforcement
  Learning (2025). \texttt{arXiv:2505.01441}
  \item Feedback-Driven Tool-Use Improvements in Large Language Models via Automated Build Environments (2025). \texttt{arXiv:2508.08791}
  \item DeepSeek-R1: Incentivizing Reasoning Capability in LLMs via Reinforcement Learning (2025). \texttt{arXiv:2501.12948}
  \item DeepResearch Bench: A Comprehensive Benchmark for Deep Research Agents (2025). \texttt{arXiv:2506.11763}
  \item Agentic Tool Use in Large Language Models (2026). \texttt{arXiv:2604.00835}
  \item Agentic Reasoning for Large Language Models (2026). \texttt{arXiv:2601.12538}
  \item The Evolution of Tool Use in LLM Agents: From Single-Tool Call to Multi-Tool Orchestration (2026). \texttt{arXiv:2603.22862}
  \item A Comprehensive Survey on Agent Skills: Taxonomy, Techniques, and Applications (2026). \texttt{arXiv:2605.07358}
  \item The Path Ahead for Agentic AI: Challenges and Opportunities (2026). \texttt{arXiv:2601.02749}
  \item Agentic Artificial Intelligence (AI): Architectures, Taxonomies, and Evaluation of Large Language Model Agents (2026). \texttt{arXiv:2601.12560}
  \item Benchmark Everything Everywhere All at Once (2026). \texttt{arXiv:2606.06462}
  \item What Twelve LLM Agent Benchmark Papers Disclose About Themselves: A Pilot Audit and an Open Scoring Schema (2026). \texttt{arXiv:2605.21404}
\end{enumerate}

\end{document}
